\clearpage\thispagestyle{empty}
\pdfbookmark[0]{Quellen zur Besprechung der Mailänder Ephemeriden für 1815}{carlini1815-sources}
\begin{center}\Large\bfseries Quellen zur Besprechung der Mailänder Ephemeriden für 1815\end{center}

Die Übertragung folgt der Besprechung von Carl Friedrich Gauß in den \emph{Werken}, Bd.~VI (Göttingen, 1874), S.~574--576. Der Erstdruck erschien in den \emph{Göttingischen gelehrten Anzeigen}, Stück~10 vom 19.~Januar 1815, S.~89--93.

Das Titelblatt von Francesco Carlinis \emph{Effemeridi astronomiche di Milano per l’anno 1815} gibt Mailand, R.\,C. Stamperia di Governo, 1814 an. Google hat das Exemplar der Harvard Science Center Library mit der Signatur Sci 295.10 digitalisiert. Es ist über Mario Carpinos digitale Bibliografie an der Sternwarte Brera des INAF zugänglich.

Barnaba Orianis Aufsatz über die Breite der Sternwarte Brera beginnt auf S.~3 des Anhangs. Auf S.~31 gibt Oriani die Breite des Turms mit $45\quad28\quad0{,}713$ an. In der entsprechenden Breitenangabe der Gaußschen Besprechung steht sowohl im Erstdruck als auch im Sammelband 49. Die Tafeln der Übertragung folgen dem Sammelband, in dem auch die Form 3c steht, wo der Erstdruck 30 hat. Orianis Aufsatz über die Marsopposition von 1813 beginnt auf S.~44.

KI-Übertragung des deutschen Originals in heutiges Deutsch.

Erstdruck (\href{https://images.p.sub.uni-goettingen.de/iiif/image/gdz:PPN319721507_1815:00000091/full/1650,/0/default.jpg}{89}, \href{https://images.p.sub.uni-goettingen.de/iiif/image/gdz:PPN319721507_1815:00000092/full/1650,/0/default.jpg}{90}, \href{https://images.p.sub.uni-goettingen.de/iiif/image/gdz:PPN319721507_1815:00000093/full/1650,/0/default.jpg}{91}, \href{https://images.p.sub.uni-goettingen.de/iiif/image/gdz:PPN319721507_1815:00000094/full/1650,/0/default.jpg}{92}, \href{https://images.p.sub.uni-goettingen.de/iiif/image/gdz:PPN319721507_1815:00000095/full/1650,/0/default.jpg}{93}).

\href{https://www1.brera.inaf.it/~mario.carpino/Biblioteca_Digitale/Effemeridi_Milano/mirror/EM1815b.pdf#page=8}{Titelblatt des Jahrbuchs}, \href{https://www1.brera.inaf.it/~mario.carpino/Biblioteca_Digitale/Effemeridi_Milano/mirror/EM1815b.pdf#page=144}{Titelblatt des Anhangs}, \href{https://www1.brera.inaf.it/~mario.carpino/Biblioteca_Digitale/Effemeridi_Milano/mirror/EM1815b.pdf#page=146}{Breite der Sternwarte Brera}, \href{https://www1.brera.inaf.it/~mario.carpino/Biblioteca_Digitale/Effemeridi_Milano/mirror/EM1815b.pdf#page=174}{Polarstern und Polhöhe}, \href{https://www1.brera.inaf.it/~mario.carpino/Biblioteca_Digitale/Effemeridi_Milano/mirror/EM1815b.pdf#page=187}{Marsopposition}, \href{https://www1.brera.inaf.it/~mario.carpino/Biblioteca_Digitale/Effemeridi_Milano/mirror/EM1815b.pdf#page=255}{Wettertabelle für Dezember 1813}.

\clearpage
